\chapter*{Introduction Générale}
\addcontentsline{toc}{chapter}{Introduction Générale}
\markboth{Introduction Générale}{Introduction Générale}

\section*{1. Contexte}
La quatrième révolution industrielle, communément désignée sous le vocable d'\textbf{Industrie 4.0}, transforme en profondeur les modes de production, d'exploitation et de maintenance des infrastructures industrielles lourdes. Cette transformation repose sur l'intégration étroite entre le monde physique (équipements électromécaniques, convoyeurs, broyeurs, capteurs industriels) et le monde numérique (systèmes informatiques distribués, traitement de flux en temps réel, algorithmique décisionnelle), donnant naissance aux \textbf{Systèmes Cyber-Physiques (Cyber-Physical Systems --- CPS)}.

Dans ce paradigme, les installations industrielles ne sont plus perçues comme de simples mécanismes isolés, mais comme des systèmes cyber-connectés générant continuellement des flux massifs de données de télémétrie (vibrations mécaniques, températures de paliers, pressions de lubrification, courants moteurs, débits volumiques). L'exploitation en temps réel de ces mesures ouvre la voie à une supervision intelligente capable d'anticiper les dérives opérationnelles et d'assurer une continuité de service optimale des installations.

Au cœur de cette dynamique, le concept de \textbf{Jumeau Numérique} (\textit{Digital Twin}), formalisé initialement par Michael Grieves \cite{grieves2014digital} puis développé par Fei Tao \cite{tao2019digital}, s'est imposé comme une innovation de rupture. Un Jumeau Numérique industriel est défini comme une réplique logicielle dynamique, synchrone et bidirectionnelle d'un équipement ou d'un processus physique. Contrairement aux outils de supervision conventionnels (SCADA) qui se limitent à afficher des valeurs brutes cloisonnées, le Jumeau Numérique contextualise la donnée :
\begin{itemize}
    \item Il intègre la \textbf{topologie structurelle} et les dépendances fonctionnelles du réseau d'équipements ;
    \item Il évalue en continu l'état de santé opérationnel des machines selon des règles métier dynamiques ;
    \item Il projette algorithmiquement l'impact d'une défaillance locale sur l'ensemble de la chaîne de valeur en aval.
\end{itemize}

\section*{2. Problématique}
Le Groupe OCP, leader mondial de l'extraction, du traitement et de la commercialisation des phosphates et dérivés phosphatés, exploite des installations caractérisées par des processus de transformation en flux continu. Une chaîne de préparation de minerai type articule successivement des unités d'extraction, de concassage primaire, de convoyage à bande, de criblage granulométrique et d'enrichissement par flottation.

Dans une telle configuration en flux continu, la survenue imprévue d'une anomalie mécanique sur une machine en amont (ex. une élévation anormale des vibrations sur un concasseur primaire) engendre inévitablement des \textbf{effets de blocage et de dégradation en cascade (effet domino)} sur l'ensemble des équipements avals. Les systèmes de supervision traditionnels actuellement en exploitation présentent deux limites majeures :
\begin{enumerate}
    \item \textbf{Absence de corrélation topologique :} Les alarmes sont émises de manière isolée au niveau de chaque automate, sans calcul automatisé de la propagation de la panne le long de la ligne de production. L'opérateur de salle de contrôle doit alors déduire mentalement et sous stress l'ensemble des machines avals menacées.
    \item \textbf{Phénomène de « tempête d'alarmes » (\textit{Alarm Flood}) :} Lorsqu'un capteur oscille autour d'une valeur seuil, les systèmes conventionnels génèrent des dizaines de notifications identiques par minute, provoquant une fatigue cognitive chez les opérateurs et augmentant le risque d'erreur humaine.
\end{enumerate}

Ces arrêts non planifiés (\textit{Unplanned Downtime}) engendrent des pertes économiques substantielles et affectent directement le Taux de Rendement Synthétique (TRS) des unités de production.

\section*{3. Objectifs}
L'objectif central de ce stage d'ingénieur en 4\textsuperscript{ème} Année Génie Informatique à l'ENSA Safi est la \textbf{conception et le développement intégral d'un prototype de Jumeau Numérique Industriel} dédié à la supervision en temps réel, à la détection automatisée d'anomalies et à l'analyse d'impact de défaillance sur une chaîne de traitement de phosphate OCP.

Les objectifs spécifiques assignés à cette réalisation logicielle sont :
\begin{itemize}
    \item \textbf{Modélisation orientée domaine (DDD) :} Formaliser les entités industrielles (Usine, Lignes, Machines, Capteurs dynamiques, Alertes, Relations de dépendance) en encapsulant les invariants métier par des gardes strictes (\textit{DomainGuard}).
    \item \textbf{Architecture logicielle découplée :} Structurer le backend selon les principes de la \textbf{Clean Architecture} et du pattern \textbf{CQRS} via le bus \textbf{MediatR}.
    \item \textbf{Moteur de traitement temps réel :} Concevoir le \textit{DigitalTwinEngine} orchestrant l'évaluation multi-seuils des anomalies et un mécanisme d'anti-rebond (\textit{AlertCooldownService}) thread-safe.
    \item \textbf{Analyse d'impact par parcours de graphe (BFS) :} Développer l'algorithme de parcours en largeur permettant de calculer dynamiquement la liste, la profondeur topologique et la sévérité estimée des machines impactées en aval.
    \item \textbf{Communication bidirectionnelle WebSockets :} Garantir la diffusion sub-seconde des changements d'état vers les clients connectés via \textbf{ASP.NET Core SignalR}.
    \item \textbf{Interface de supervision réactive sous Angular 19 :} Développer un tableau de bord moderne exploitant les \textit{Signals}, une cartographie topologique interactive sous \textit{vis-network} et des graphiques de télémétrie sous \textit{Chart.js}.
    \item \textbf{Validation expérimentale et tests automatisés :} Assurer une couverture complète par 83 tests unitaires \textit{xUnit} et valider le comportement du système sur un scénario déterministe de panne de concasseur.
\end{itemize}

\section*{4. Méthodologie}
Afin de concilier la rigueur d'ingénierie logicielle et les contraintes industrielles d'OCP, la démarche méthodologique adoptée s'articule autour de :
\begin{itemize}
    \item \textbf{La démarche Agile / Scrum :} Découpage du projet en sprints de développement avec revues régulières pour valider l'adéquation des fonctionnalités avec les besoins de supervision.
    \item \textbf{L'approche Domain-Driven Design :} Établissement d'un langage ubiquitaire (\textit{Ubiquitous Language}) pour modéliser fidèlement les équipements (roue-pelle, concasseur, crible, cellule de flottation) et leurs interconnexions.
    \item \textbf{Le développement dirigé par les tests (TDD) :} Écriture préalable des tests unitaires pour valider les règles métier, les transitions d'état et les algorithmes de parcours de graphe.
    \item \textbf{Le prototypage expérimental autonome :} Implémentation d'un \textit{Worker Service} de simulation capable de générer des signaux physiques réalistes et de rejouer des scénarios d'avarie sans risque pour les installations réelles.
\end{itemize}

\section*{5. Organisation du Rapport}
Le présent document est structuré en sept chapitres thématiques :
\begin{itemize}
    \item Le \textbf{Chapitre 1} présente l'organisme d'accueil (Groupe OCP), ses activités, son site et l'entité d'accueil.
    \item Le \textbf{Chapitre 2} analyse les principes de supervision industrielle, pose la problématique et formalise le cahier des charges.
    \item Le \textbf{Chapitre 3} détaille la conception architecturale, le modèle de données, les cas d'utilisation, les diagrammes de séquence et la machine d'états.
    \item Le \textbf{Chapitre 4} décrit la réalisation technique du backend .NET 9, de la simulation, de la détection d'anomalies, du hub SignalR et de l'interface Angular 19.
    \item Le \textbf{Chapitre 5} formalise la modélisation sous forme de graphe, détaille l'algorithme BFS d'analyse d'impact et présente l'exécution du scénario de défaillance du concasseur CR-001.
    \item Le \textbf{Chapitre 6} expose la stratégie de test, la suite des 83 tests unitaires, les mesures de latence et la discussion des résultats.
    \item Le \textbf{Chapitre 7} identifie les limites actuelles et trace les perspectives d'évolution industrielle.
    \item La \textbf{Conclusion Générale} synthétise les apports techniques et personnels du stage.
\end{itemize}